\begin{table}[!htbp]
\centering
\scriptsize
\setlength{\tabcolsep}{3pt}
\renewcommand{\arraystretch}{1.16}
\caption{Scenario-specific model-selection guidance. Recommendations are conditional and do not turn any baseline into a proposed algorithm.}
\label{tab:model-selection-guidance}
\begin{tabularx}{\linewidth}{@{}>{\RaggedRight\arraybackslash}p{0.14\linewidth}>{\RaggedRight\arraybackslash}p{0.29\linewidth}>{\RaggedRight\arraybackslash}X@{}}
\toprule
Scenario & Recommended strategy & Main caution \\
\midrule
label-rich & Framework-Balanced or Fixed-LightGBM & Do not claim broad dominance; differences are small. \\
label-scarce & Fixed-LightGBM with validation-only threshold; Framework-Balanced when multiple candidates are available & Normal-only and 5\% settings remain weak; avoid overclaiming. \\
safety-critical & Framework-Safety & FAR increases substantially, so alarm workflow must tolerate more warnings. \\
low false alarm & Framework-Low-False-Alarm & MDR increases; not suitable where missed detections dominate. \\
domain shift & Fixed-LightGBM or cost-specific framework strategy & Framework does not uniformly beat LightGBM under domain shift. \\
missing sensor & Framework-Robust or Framework-Balanced & Needs deployment-time validation for exact missing-sensor pattern. \\
low-latency candidate selection & Fixed-LightGBM or Fixed-XGBoost & Feature extraction cost must be included in real deployment. \\
\bottomrule
\end{tabularx}
\end{table}
